\subsection{李群情形}

命题 11. 假设 k 是 R 、C 或一个非离散完备超度量域。设 h 是 g 的一个分裂嘉当子代数。

\begin{enumerate}
\def\labelenumi{(\roman{enumi})}
\item
  \(\operatorname{Aut}(\mathfrak{g},\mathfrak{h})\) 是 \(\operatorname{Aut}(\mathfrak{g})\) 的一个李子群，其李代数是 ad h。
\item
  \(f(\mathrm{T}_{\mathrm{Q}})\) 与 \((q\circ f)(\mathrm{T}_{\mathrm{P}})\) 是 \(\operatorname {Aut}(\mathfrak{g},\mathfrak{h})\) 的开子群
\item
  \(\mathrm{Aut}_e(\mathfrak{g})\) 是 \(\mathrm{Aut}(\mathfrak{g})\) 的一个开子群。
\item
  若 \(k = \mathbf{R}\) 或 \(\mathbf{C}\) ，则 \(\operatorname{Aut}_e(\mathfrak{g})\) 是 \(\operatorname{Aut}(\mathfrak{g})\) 的单位元连通分支，换言之即 \(\operatorname{Int}(\mathfrak{g})\) 。
\end{enumerate}

由第三章 §3 第 8 节 命题 29 的推论 2 与第 10 节 命题 36，\(\operatorname{Aut}(\mathfrak{g}, \mathfrak{h})\) 是 \(\operatorname{Aut}(\mathfrak{g})\) 的李子群，其李代数是这样的 \(\operatorname{ad} x (x \in \mathfrak{g})\) 组成的集合：使 \((\operatorname{ad} x) \mathfrak{h} \subset \mathfrak{h}\) ，换言之即 \(\operatorname{ad} \mathfrak{h}\) 。

设 \(\mathrm{H} \in \mathfrak{h}\) 。存在 \(\varepsilon > 0\) 具有下列性质：对 \(t \in k\) 与 \(|t| < \varepsilon\) ，\(\exp(t\gamma(\mathrm{H}))\) 对一切 \(\gamma \in \mathrm{P}(\mathrm{R})\) 有定义，且映射 \(\gamma \mapsto \exp(t\gamma(\mathrm{H}))\) 是同态 \(\sigma_t\) （从 \(\mathrm{P}(\mathrm{R})\) 到 \(k^*\) 的）。对 \(|t| < \varepsilon\) ，\(\exp(t\operatorname{ad}\mathrm{H})\) 有定义，在 \(\mathfrak{h}\) 上诱导恒等，且在 \(\mathfrak{g}^\alpha\) 上诱导比为 \(\sigma_t(\alpha)\) 的位似；故 \(\exp t\operatorname{ad}\mathrm{H} \in (q \circ f)(\mathrm{T}_{\mathrm{P}})\) 。这就证明，鉴于 (i)，\((q \circ f)(\mathrm{T}_{\mathrm{P}})\) 包含 \(\operatorname{Aut}(\mathfrak{g},\mathfrak{h})\) 中 1 的一个邻域，从而是 \(\operatorname{Aut}(\mathfrak{g},\mathfrak{h})\) 的开子群。更有甚者，\(f(\mathrm{T}_{\mathrm{Q}})\) 是 \(\operatorname{Aut}(\mathfrak{g},\mathfrak{h})\) 的开子群。

对一切 \(\alpha \in \mathrm{R}\) ，\(\exp \operatorname{ad} \mathfrak{g}^{\alpha} \subset \operatorname{Aut}_{e}(\mathfrak{g})\) 。鉴于 (ii)，\(\operatorname{Aut}_{e}(\mathfrak{g})\) 包含 \(\operatorname{Aut}(\mathfrak{g})\) 中 1 的一个邻域，这就证明了 (iii)。

假设 \(k = \mathbf{R}\) 或 \(\mathbf{C}\) 。则 \(\mathrm{Aut}_e(\mathfrak{g})\) 包含在单位元连通分支 \(\mathrm{C}\) （\(\mathrm{Aut}(\mathfrak{g})\) 的）中（第七章 §3 第 1 节），并由 (iii) 知在 \(\mathrm{Aut}(\mathfrak{g})\) 中是开的。于是 \(\mathrm{Aut}_e(\mathfrak{g}) = \mathrm{C}\) 。最后，由第三章 §9 第 8 节 命题 30 (i)，\(\mathrm{C} = \mathrm{Int}(\mathfrak{g})\) 。

